\documentclass[a4paper,10pt]{article}
\usepackage[ngerman]{babel}
\usepackage[utf8]{inputenc}
\usepackage{wasysym}% provides \ocircle and \Box
\usepackage{enumitem}% easy control of topsep and leftmargin for lists
\usepackage{color}% used for background color
\usepackage{forloop}% used for \Qrating and \Qlines
\usepackage{ifthen}% used for \Qitem and \QItem
\usepackage{typearea}
\usepackage{array}
\usepackage{hyperref}
\areaset{17cm}{26cm}
\setlength{\topmargin}{-1cm}
\usepackage{scrlayer-scrpage}
\pagestyle{scrheadings}
\ihead{Developers' and Researchers' Perspective on Zero-Knowledge Security}
\ohead{\pagemark}
\chead{}
\cfoot{}

%%%%%%%%%%%%%%%%%%%%%%%%%%%%%%%%%%%%%%%%%%%%%%%%%%%%%%%%%%%%
%% Beginning of questionnaire command definitions %%
%%%%%%%%%%%%%%%%%%%%%%%%%%%%%%%%%%%%%%%%%%%%%%%%%%%%%%%%%%%%
%%
%% 2010, 2012 by Sven Hartenstein
%% mail@svenhartenstein.de
%% http://www.svenhartenstein.de

%% \Qq = Questionaire question. Oh, this is just too simple. It helps
%% making it easy to globally change the appearance of questions.
\newcommand{\Qq}[1]{\textbf{#1}}

%% \QO = Circle or box to be ticked. Used both by direct call and by
%% \Qrating and \Qlist.
\newcommand{\QO}{$\Box$}% or: $\ocircle$


%% \Qline = Again, this is very simple. It helps setting the line
%% thickness globally. Used both by direct call and by \Qlines.
\newcommand{\Qline}[1]{\noindent\rule{#1}{0.6pt}}

%% \Qlines = Insert NUM lines with width=\linewith. You can change the
%% \vskip value to adjust the spacing.
\newcounter{ql}
\newcommand{\Qlines}[1]{\forloop{ql}{0}{\value{ql}<#1}{\vskip0em\Qline{\linewidth}}}

%% \Qlist = This is an environment very similar to itemize but with
%% \QO in front of each list item. Useful for classical multiple
%% choice. Change leftmargin and topsep accourding to your taste.
\newenvironment{Qlist}{%
\renewcommand{\labelitemi}{\QO}
\begin{itemize}[leftmargin=1.5em,topsep=-.5em]
}{%
\end{itemize}
}

%% \Qitem = Item with automatic numbering. 
\newcounter{itemnummer}
\newcommand{\Qitem}[2][]{% #1 optional, #2 notwendig
\ifthenelse{\equal{#1}{}}{\stepcounter{itemnummer}}{}
\ifthenelse{\equal{#1}{a}}{\stepcounter{itemnummer}}{}
\begin{enumerate}[topsep=2pt,leftmargin=2.8em]
\item[\textbf{\arabic{itemnummer}#1.}] #2
\end{enumerate}
}

%% \QItem = Like \Qitem but with alternating background color. This
%% might be error prone as I hard-coded some lengths (-5.25pt and
%% -3pt)! I do not yet understand why I need them.
\definecolor{bgodd}{rgb}{0.8,0.8,0.8}
\definecolor{bgeven}{rgb}{0.9,0.9,0.9}
\newcounter{itemoddeven}
\newlength{\gb}
\newcommand{\QItem}[2][]{% #1 optional, #2 notwendig
\setlength{\gb}{\linewidth}
\addtolength{\gb}{-5.25pt}
\ifthenelse{\equal{\value{itemoddeven}}{0}}{%
\noindent\colorbox{bgeven}{\hskip-3pt\begin{minipage}{\gb}\Qitem[#1]{#2}\end{minipage}}%
\stepcounter{itemoddeven}%
}{%
\noindent\colorbox{bgodd}{\hskip-3pt\begin{minipage}{\gb}\Qitem[#1]{#2}\end{minipage}}%
\setcounter{itemoddeven}{0}%
}
}

%%%%%%%%%%%%%%%%%%%%%%%%%%%%%%%%%%%%%%%%%%%%%%%%%%%%%%%%%%%%
%% End of questionnaire command definitions %%
%%%%%%%%%%%%%%%%%%%%%%%%%%%%%%%%%%%%%%%%%%%%%%%%%%%%%%%%%%%%
\begin{document}

% ##########################################################
\begin{center}
    \textbf{\huge Developers' and Researchers' Perspective on Zero-Knowledge Security}
\end{center}\vskip1em
% ##########################################################

% ##########################################################
\section*{Introduction}
This questionnaire is part of a research project conducted by Stefanos Chaliasos, Alex Hicks, Arman Kolozyan, and Tom Sorger, studying the current state of zero-knowledge security from the perspective of developers and researchers building ZK systems.
Your participation is \textbf{anonymous}.
To preserve participants' privacy, only aggregated results will be shared publicly: as part of a research paper, a blog post, and a dedicated page on \href{https://zkbugs.com/}{zkbugs}.
Individual responses will not be published.
The questionnaire should take approximately \textbf{5--10 minutes} to complete.
Feel free to \textit{skip any question that you would prefer not to answer}, or that does not apply to you.
Thank you for contributing your time and expertise.
% ##########################################################

% ##########################################################
\section{Demographics}
% The following questions are intended to gauge participants' experience level and to understand the nature of the ZK systems and protocols they work on.

% Q1
\QItem{ \Qq{Please specify your software development experience in years.}
\begin{Qlist}
    \item [A] 0--2
    \item [B] 3--5
    \item [C] 6--10
    \item [D] 11+
\end{Qlist}
}

% Q2
\QItem{ \Qq{Please specify your ZKP development experience in years.}
\begin{Qlist}
    \item [A] Less than 1
    \item [B] 1--2
    \item [C] 3--5
    \item [D] 6+
\end{Qlist}
}

% Q3
\QItem{ \Qq{What is the size of the engineering (ZK-focused) team in the organisation you work for?}
\begin{Qlist}
    \item [A] 1--5
    \item [B] 6--25
    \item [C] 26+
    \item [D] I'm an independent auditor
\end{Qlist}
}

% Q4
\QItem{ \Qq{What best describes the deployment stage of the most advanced ZK project you have contributed to?}
\begin{Qlist}
\item [A] Design/specification only
\item [B] Academic/research prototype (not production-bound)
\item [C] Active development
\item [D] Public testnet or staging
\item [E] Mainnet or production deployment
\end{Qlist}
}

% Q5
\QItem{ \Qq{Which ZKP system do you work with? Select all that apply.}
\begin{Qlist}
\item [A] Circuit DSLs (e.g., Circom, Noir, Cairo, Leo, Halo2, Plonky3)
\item [B] zkVMs (e.g., RISC Zero, SP1, Jolt, zkWasm)
\item [C] zkEVMs (e.g., Scroll, Polygon zkEVM)
\item [D] zkVM Guest (e.g., Bridges implemented with zkVMs)
\item [E] Other: \Qline{13.3cm}
\end{Qlist}
}

% Q6
\QItem{ \Qq{In which primary domain(s) are the ZK systems you develop deployed?}
\begin{Qlist}
\item [A] Public Blockchain Infrastructure (e.g., L1/L2 scaling, zkEVMs, Bridges)
\item [B] Decentralized Applications (dApps) (e.g., DeFi, Privacy Mixers, On-chain Gaming)
\item [C] Digital Identity \& Credentials (e.g., zk-ID, KYC/AML compliance, Sybil resistance)
\item [D] Privacy Infrastructure (e.g., privacy enabling L1/L2s)
\item [E] Verifiable Computing \& AI (e.g., ZK-ML, verifying outsourced computation)
\item [F] Libraries/Infra (e.g., cryptographic libraries)
\item [G] Other: \Qline{13cm}
\end{Qlist}
}

% Q7
\QItem{ \Qq{What layer of the ZKP stack do you primarily work with? Select all that apply.}
\begin{Qlist}
\item [A] Circuit/constraint correctness
\item [B] VM semantics 
\item [C] Proof system 
\item [D] Verifiers
\item [E] Compilation Pipelines
\item [F] Other: \Qline{13.3cm}
\end{Qlist}
}

% Q8
\QItem{ \Qq{Which resources do you use to get informed about ZKP security and the latest exploits? Select all that apply.}
\begin{Qlist}
    \item [A] Research papers (e.g., ePrint, conference proceedings)
    \item [B] Social media (e.g., X (Twitter), YouTube)
    \item [C] Blog posts (e.g., Medium)
    \item [D] Security company reports (e.g., zkSecurity, Veridise, Trail of Bits)
    \item [E] ZKP community forums (e.g., ZKProof, Ethereum Research)
    \item [F] Colleagues
    \item [G] Conferences/workshops
    \item [H] Documentation
    \item [I] GitHub repositories (e.g., bug trackers, issues)
    \item [J] Other: \Qline{13.3cm}
\end{Qlist}
}

% Q9
\QItem{ \Qq{Regarding the resources selected in the previous question, how would you characterize the current rate of new security information being produced?}
\begin{Qlist}
    \item [A] High Velocity: New vulnerabilities and techniques emerge weekly; I must check these resources constantly to stay relevant.
    \item [B] Steady Growth: There is a consistent stream of new information, but the core security principles remain largely unchanged month-to-month.
    \item [C] Static/Slow: The ``classic'' papers and docs still cover 90\% of what I need; new exploits are rare or mostly variations of known issues.
    \item [D] Fragmented: Information is being produced, but it is too scattered across resources to track effectively.
\end{Qlist}
}

% Q10
\QItem{ \Qq{How knowledgeable are you about zero-knowledge security from 1 (not familiar) to 5 (expert)?}
\begin{Qlist}
\item [1] Not Familiar
\item [2] 
\item [3] 
\item [4]
\item [5] Expert
\end{Qlist}
}
% ##########################################################

\section{Security Effort \& AI Impact }

% Q11
\QItem{ \Qq{What is the average duration of a full external ZK security audit for projects you work on?}
\begin{Qlist}
    \item [A] One week or less
    \item [B] 2--4 weeks
    \item [C] 1--2 months
    \item [D] 3 months+
\end{Qlist}
}

% Q12
\QItem{ \Qq{How much effort is typically required to integrate security analysis or formal verification tools into your ZK codebase (e.g., specifications, frontends, CI)?}
\begin{Qlist}
    \item [A] Minimal (Plug-and-play)
    \item [B] Low (Few hours of configuration)
    \item [C] Moderate (1-3 days of custom coding/specs)
    \item [D] High (Requires a dedicated sub-phase of the audit)
\end{Qlist}
}

% Q13
\QItem{ \Qq{How has the introduction of AI/LLMs affected your ZK development workflow or team productivity?}
\begin{Qlist}
    \item [A] Reduced costs/time significantly (more than 25\% efficiency gain)
    \item [B] Marginal efficiency gains, but delivery expectations from stakeholders remain the same
    \item [C] No impact on timing or workload yet
    \item [D] Increased time/cost due to manual verification of AI-generated code or suggestions
    \item [E] Other: \Qline{13.3cm}
\end{Qlist}
}

% Q14
\QItem{ \Qq{How has the introduction of AI/LLMs affected the security of your ZK codebases during development?}
\begin{Qlist}
    \item [A] Less secure (i.e., bugs introduced by the AI and went uncovered during development)
    \item [B] More secure (i.e., AI detected bugs during development)
    \item [C] No impact
    \item [D] Other: \Qline{13.3cm}
\end{Qlist}
}

% ##########################################################

\section{Familiarity/experience with security tools}

% Q15
\QItem{ \Qq{Which of the following types of tools have you used for ZKP security analysis? Including potentially internal or paid tools. Select all that apply.}
\begin{Qlist}
    \item [A] Static analyser (e.g., Circomspect, Custom Static Analysis Tools)
    \item [B] Theorem Provers/DSLs (e.g., Clean, zkLean, Lean) 
    \item [C] Automated (or semi-automated) Formal Verification (e.g., Picus, Civer)
    \item [D] Fuzzer for circuits or ZKVMs (e.g., zkFuzz, Arguzz)
    \item [E] LLMs (Interactive sessions with Claude/Codex with or without agent skills)
    \item [F] Customised Property-based testing
    \item [G] Automated tools leveraging AI (e.g., zkao)
    \item [H] Other: \Qline{13.3cm}
\end{Qlist}
}

% Q16
\QItem{ \Qq{Name up to five tools you frequently use. Write internal + tool type for non-public tools (e.g., Internal + Static Analysis)}:\vspace{0.4cm}\\
\vspace{0.4cm}\Qline{15cm}\\ 
\vspace{0.4cm}\Qline{15cm}\\ 
\vspace{0.4cm}\Qline{15cm}\\ 
\vspace{0.4cm}\Qline{15cm}\\ 
\Qline{15cm}
}

% Q17
\QItem{ \Qq{Below we enumerate a number of security analysis characteristics. To what extent do you agree regarding their importance?}

\begin{tabular}{|>{\centering\arraybackslash}p{3cm}|*{5}{>{\centering\arraybackslash}p{2cm}|}}
    \hline
& \textbf{Strongly Disagree} & \textbf{Disagree} & \textbf{Neutral} & \textbf{Agree} & \textbf{Strongly Agree} \\
\hline
Low false positives (tool does not wrongly indicate that there is a vulnerability) & \QO & \QO & \QO & \QO & \QO \\
\hline
Ease of use (e.g., simple CLI, integrated into IDE) & \QO & \QO & \QO & \QO & \QO \\
\hline
Documentation & \QO & \QO & \QO & \QO & \QO \\
\hline
Report quality (i.e., detail of vulnerabilities detected) & \QO & \QO & \QO & \QO & \QO \\
\hline
Scalability/performance (Time to report) & \QO & \QO & \QO & \QO & \QO \\
\hline
Producing PoC & \QO & \QO & \QO & \QO & \QO \\
\hline
Cost (Free vs Paid options) & \QO & \QO & \QO & \QO & \QO \\
\hline
Integration with CI/CD pipelines & \QO & \QO & \QO & \QO & \QO \\
\hline
Confidence of completeness & \QO & \QO & \QO & \QO & \QO \\
\hline
\end{tabular}
}

% Q18
\QItem{ \Qq{What other factors could positively or negatively affect you when using a security analysis tool?}\vspace{0.4cm}\\
\vspace{0.4cm}\Qline{15cm}\\ 
\vspace{0.4cm}\Qline{15cm}\\ 
\vspace{0.4cm}\Qline{15cm}\\ 
\vspace{0.4cm}\Qline{15cm}\\ 
\Qline{15cm}
}

% Q19
\QItem{ \Qq{Which vulnerabilities do you think are more difficult to detect during development?}
\begin{Qlist}
    \item [A] Under-constrained circuits (missing constraints allowing unsound proofs)
    \item [B] Over-constrained circuits (constraints that are too strict, blocking valid witnesses)
    \item [C] Arithmetic field errors (e.g., unintended modular reduction, non-native field arithmetic)
    \item [D] Bit decomposition, range-check, or overflow/underflow errors
    \item [E] Incorrect public input handling or signal aliasing
    \item [F] Unsound proof system implementation or configuration
    \item [G] Logic/semantic errors (circuit does not match intended specification or constraint and witness divergences)
    \item [H] Hash function or cryptographic primitive misuse in circuits
    \item [I] Fiat-Shamir related attacks
    \item [J] Incorrect integration in a system
    \item [K] Other: \Qline{13.3cm}
\end{Qlist}
}

% Q20
\QItem{ \Qq{Which vulnerabilities do you think are more adequately covered by security tools?}
\begin{Qlist}
    \item [A] Under-constrained circuits (missing constraints allowing unsound proofs)
    \item [B] Over-constrained circuits (constraints that are too strict, blocking valid witnesses)
    \item [C] Arithmetic field errors (e.g., unintended modular reduction, non-native field arithmetic)
    \item [D] Bit decomposition, range-check, or overflow/underflow errors
    \item [E] Incorrect public input handling or signal aliasing
    \item [F] Unsound proof system implementation or configuration
    \item [G] Logic/semantic errors (circuit does not match intended specification or constraint and witness divergences)
    \item [H] Hash function or cryptographic primitive misuse in circuits
    \item [I] Fiat-Shamir related attacks
    \item [J] Incorrect integration in a system
    \item [K] Other: \Qline{13.3cm}
\end{Qlist}
}

% Q21
\QItem{ \Qq{Which vulnerabilities that you deem crucial \textbf{cannot} be detected by \textbf{automated security} tools (including LLMs) that you have used? Select all that apply.}
\begin{Qlist}
    \item [A] Under-constrained circuits (missing constraints allowing unsound proofs)
    \item [B] Over-constrained circuits (constraints that are too strict, blocking valid witnesses)
    \item [C] Arithmetic field errors (e.g., unintended modular reduction, non-native field arithmetic)
    \item [D] Bit decomposition, range-check, or overflow/underflow errors
    \item [E] Incorrect public input handling or signal aliasing
    \item [F] Unsound proof system implementation or configuration
    \item [G] Logic/semantic errors (circuit does not match intended specification or constraint and witness divergences)
    \item [H] Hash function or cryptographic primitive misuse in circuits
    \item [I] Fiat-Shamir related attacks
    \item [J] Incorrect integration in a system
    \item [K] Other: \Qline{13.3cm}
\end{Qlist}
}

% Q22
\QItem{ \Qq{Do you find security tools helpful when developing or reviewing ZK circuits and systems?}
\begin{Qlist}
\item [1] Counter-productive
\item [2] 
\item [3] 
\item [4]
\item [5] Essential
\end{Qlist}
}

% Q23
\QItem{ \Qq{For which languages or frameworks do you use security tools (\textbf{excluding interactive LLMs})? Select all that apply.}
\begin{Qlist}
    \item [A] Circom
    \item [B] Halo2 
    \item [C] Zirgen
    \item [D] Gnark
    \item [E] Plonky3
    \item [F] zkVM Guest Programs
    \item [G] Noir
    \item [H] Leo
    \item [I] Cairo
    \item [J] Other: \Qline{13.3cm}
\end{Qlist}
}

% Q24
\QItem{ \Qq{What percentage of security researchers do you believe AI auditing currently outperforms?}
\begin{Qlist}
    \item [A] 0\%
    \item [B] 1--25\%
    \item [C] 26--50\%
    \item [D] 51--75\%
    \item [E] 76--99\%
    \item [F] 100\%
\end{Qlist}
}

% Q25
\QItem{ \Qq{What percentage of security researchers do you believe AI auditing will outperform in the next 3 years?}
\begin{Qlist}
    \item [A] 0\%
    \item [B] 1--25\%
    \item [C] 26--50\%
    \item [D] 51--75\%
    \item [E] 76--99\%
    \item [F] 100\%
\end{Qlist}
}

% Q26
\QItem{ \Qq{Which tools do you think are mostly missing on the market, and would you use if available?}
\begin{Qlist}
    \item [A] Quick Static Analysis or Linters
    \item [B] Automated Formal Verification
    \item [C] Domain Specific Interactive Provers (with LLM support)
    \item [D] LLM Skills
    \item [E] AI Tools
    \item [F] Fuzzers
    \item [G] Proof Assistants
\end{Qlist}
}

% Q27
\QItem{ \Qq{Have you ever tried a ZK security tool and decided not to continue using it? If so, why?}\vspace{0.4cm}\\
\vspace{0.4cm}\Qline{15cm}\\ 
\vspace{0.4cm}\Qline{15cm}\\ 
\vspace{0.4cm}\Qline{15cm}\\ 
\vspace{0.4cm}\Qline{15cm}\\ 
\Qline{15cm}
}

% Q28
\QItem{ \Qq{In which tasks do you think LLMs are extremely good already? Select all that apply.}
\begin{Qlist}
    \vspace{0.4cm}
    \item [A] Nothing
    \item [B] Validating bugs and producing PoC
    \item [C] Finding bugs in Circuit Code
    \item [D] Finding bugs in non-ZK application code (e.g., Rust)
    \item [E] Finding bugs in proof system code
    \item [F] Finding bugs in zkVM guest programs
    \item [G] Everything
\end{Qlist}
}

% Q29
\QItem{ \Qq{When using Formal Verification (FV) or automated tools, how concerned are you about the Trusted Computing Base (TCB) (e.g., bugs in the solver, compiler, or the tool itself)?}
\begin{Qlist}
    \item [A] Not concerned: The risk is negligible compared to human error.
    \item [B] Somewhat concerned: We perform sanity checks on tool outputs.
    \item [C] Highly concerned: We use multiple tools or manual cross-checks to mitigate TCB risks.
\end{Qlist}
}

% ##########################################################

\section{Security practices you apply (or applied by other team members in your current organisation)}
If you are an independent developer or researcher, answer the following questions on how you apply security practices.

% Q30
\QItem{ \Qq{What techniques does the organisation you work for employ when developing and securing ZK circuits and systems?}
\begin{Qlist}
    \item [A] Manual verification
    \item [B] Formal verification 
    \item [C] Differential testing (comparing circuit output against a reference implementation)
    \item [D] Property-based or fuzz testing
    \item [E] Running/extending developer test suites (e.g., unit and integration tests)
    \item [F] Using open source security tools
    \item [G] Using internal security tools
    \item [H] Using LLMs interactively
    \item [I] Using customized LLM Skills
    \item [J] Using automated AI tools
    \item [K] Other: \Qline{13.3cm}
\end{Qlist}
}

% Q31
\QItem{ \Qq{What are the security practices that your organisation follows? Select all that apply.}
\begin{Qlist}
    \item [A] Code Reviews
    \item [B] Internal Auditing
    \item [C] External Auditing
    \item [D] Bug Bounty 
    \item [E] Runtime monitoring 
    \item [F] Testing (e.g., unit testing, integration testing) 
    \item [G] Other: \Qline{13.3cm}
\end{Qlist}
}

% Q32
\QItem{ \Qq{What kind of security tools are you using inside the organisation you work for? Select all that apply.}
\begin{Qlist}
    \item [A] Open-source
    \item [B] Internally developed
    \item [C] Extended existing open-source tools
    \item [D] Third-party service
    \item [E] We don't use any security tools
\end{Qlist}
}

% Q33
\QItem{ \Qq{Does the organisation you work for develop in-house security tools? Select all that apply.}
\begin{Qlist}
    \item [A] AI-Assisted Decision Support
    \item [B] Autonomous AI Security Systems
    \item [C] Traditional (Non-AI) Security Tools
    \item [D] No In-House Development
    \item [E] Other: \Qline{13.3cm}
\end{Qlist}
}

% Q34
\QItem{ \Qq{If the organisation you work for runs any tool that requires specifications (e.g., for formal verification of circuit correctness), who writes them?}
\begin{Qlist}
    \item [A] Circuit/protocol developers
    \item [B] Auditors
    \item [C] Auditors and developers collaboratively
    \item [D] I don't know
    \item [E] Other: \Qline{13.3cm}
\end{Qlist}
}

% Q35
\QItem{ \Qq{How are you validating their specifications?}\vspace{0.4cm}\\
\vspace{0.4cm}\Qline{15cm}\\ 
\vspace{0.4cm}\Qline{15cm}\\ 
\vspace{0.4cm}\Qline{15cm}\\ 
\vspace{0.4cm}\Qline{15cm}\\ 
\Qline{15cm}
}

% Q36
\QItem{ \Qq{What is the most common reason your organisation might push back against adopting certain automated or formal verification tools during development?}
\begin{Qlist}
    \item [A] Cost/Extra budget requirements
    \item [B] Timeline: Belief that it will slow down delivery
    \item [C] Skepticism: Lack of trust in ``mathematical proof'' over manual review
    \item [D] Other: \Qline{13.3cm}
\end{Qlist}
}

% Q37
\QItem{ \Qq{To what extent do you agree with each of the following statements?}

\begin{tabular}{|>{\centering\arraybackslash}p{3cm}|*{5}{>{\centering\arraybackslash}p{2cm}|}}
    \hline
& \textbf{Strongly Disagree} & \textbf{Disagree} & \textbf{Neutral} & \textbf{Agree} & \textbf{Strongly Agree} \\
\hline
The feedback of security audits improves the quality of the audited application & \QO & \QO & \QO & \QO & \QO \\
\hline
Most of the issues reported on security audits are resolved before shipping an application & \QO & \QO & \QO & \QO & \QO \\
\hline
A primary reason to do a security audit is that they are typically considered important & \QO & \QO & \QO & \QO & \QO \\
\hline
\end{tabular}
}

% Q38
\QItem{ \Qq{Is there anything else you would like to add?}\vspace{0.4cm}\\
\vspace{0.4cm}\Qline{15cm}\\ 
\vspace{0.4cm}\Qline{15cm}\\ 
\vspace{0.4cm}\Qline{15cm}\\ 
\vspace{0.4cm}\Qline{15cm}\\ 
\Qline{15cm}
}
% ##########################################################
\end{document}
